%
% 5_Fortsetzung.tex -- Abschnitt 5 Fortsetzende Idee der Resultanten Berechnung
%
% (c) Jakob Gierer, OST Ostschweizer Fachhochschule
%
% !TEX root = ../../buch.tex
% !TEX encoding = UTF-8
%
\section{Geht das noch besser?
\label{resultante:section:fortsetzung}}
\kopfrechts{Polynom-Restfolgen}
Der euklidische Algorithmus zur Berechnung grösster gemeinsamer Teiler
zweier ganzer Zahlen ist mehr als 2000 Jahre alt.
\index{euklidischer Algorithmus}%
Es ist der älteste bekannte Algorithmus.
Beim Versuch, den euklidischen Algorithmus zu verallgemeinern, um den
grössten gemeinsamen Teiler oder die Resultante von zwei Polynomen
zu berechnen, erreicht man schnell das Konzept der Polynom-Restfolge
(polynomial remainder sequence, PRS).
Dies führt zu schnell anwachsenden Koeffizienten.
Dieses schnelle Anwachsen ist zum Glück nicht Teil des Problems, sondern
ein Artefakt der naiven Ansätze der Verallgemeinerung.
Das Wachstum lässt sich kontrollieren, indem von jedem Polynom der PRS
ein nicht benötigter Teil entfernt wird.
Was dieser Teil ist und wie er entfernt wird, erschliesst sich im
folgenden Abschnitt.

Der Schlüssel, um dieses Koeffizienten-Wachstum zu kontrollieren,
ohne die aufwendige Berechunug des \(\ggt\) der Koeffizienten, ist die
Grundlage der Subresultanten.
Das Theorem von Subresultanten aus \cite{resultant:PRSalg} zeigt, dass
\index{Subresultante}%
jedes Polynom einer PRS proportional zu einer Subresultanten der ersten
zwei Polynome dieser Folge ist.


\subsection{Grundlegender Aufbau
\label{resultante:subsection:aufbau}}
Für die Konstruktion einer PRS wird die Division mit Rest benötigt,
dabei kann dies beliebig normiert werden.
Der euklidische Algorithmus teilt zwei Polynome in \(\mathbb{Q}\),
wobei das Polynom mit dem grösseren Grad der Dividend ist und das mit
dem kleineren der Divisor:
\begin{equation}
    P_i / P_{i+1} = Q_i \quad \text{mit}\quad R_i.
\end{equation}
Als nächster Schritt wird der Rest \(R_i\) der Division zum neuen Divisor
\(P_{i+2}\) und der alte Divisor wird zum neuen Dividend \(P_{i+1}\).
Durch die Division können problematische Brüche als Koeffizienten entstehen.
\begin{beispiel}
    Mit den Polynomen \(a(X)=4X^4-7X^3+22X^2-24X+5\) und
    \(b(X)=8X^3+8X^2-16X\) soll der euklidische Algorithmus gezeigt
    werden.
    Begonnen wird mit \(P_1=a(X)\) und \(P_2=b(X)\), da
    \(\deg(a)>\deg(b)\) ist, dabei wird
    \begin{equation}
        (4X^4-7X^3+22X^2-24X+5)/(8X^3+8X^2-16X).
    \end{equation}
    Nach der Polynom-Division erhählt man
    \(Q_1=\frac{X}{2}-\frac{11}{8}\) und \(R_1=P_3=41X^2-46X+5\).
    Die nächste Division ist
    \begin{equation*}
        P_2/P_3,
    \end{equation*}
    daraus folgt die neue Division
    \begin{equation}
        (8X^3+8X^2-16X)/(41X^2-46X+5),
    \end{equation}
    was \(R_2=P_4=\frac{3480}{1681}X-\frac{3480}{1681}\) ergibt.
    Da \(R_2\neq0\) ist, geht der Algorithmus weiter mit
    \begin{equation*}
        P_3/P_4,
    \end{equation*}
    und
    \begin{equation}
        (41X^2-46X+5)\Big/\biggl(\frac{3480}{1681}X-\frac{3480}{1681}\biggr),
    \end{equation}
    daraus resultiert \(R_3=P_5=0\).
    \label{resultante:beispiel:eualg}
\end{beispiel}
Wie man in Beispiel \ref{resultante:beispiel:eualg} sieht, entstehen
bei der Division grosse Brüche.
Zudem ist der Grad jedes weiteren Polynoms um eins kleiner als der
des vorherigen.

\subsection{Verbesserungsmöglichkeiten
\label{resultante:subsection:besser}}
Der euklidische Algorithmus liefert eine Folge von Quotienten und Resten.
Diese Quotienten und Reste führen in der Standardvorgehensweise zu
Brüchen, deren Zähler und Nenner riesig werden können.
Dies ist der Grund für die verallgemeinerte PRS, mit den folgenden
Zielen zur Optimierung:
\begin{enumerate}[label=\alph*)]
    \item Rechnen nur mit ganzen Zahlen, da dies effizienter ist (kein
    Kürzen von Brüchen).
    \item Nur mit kleinen Koeffizienten rechnen.
\end{enumerate}

\subsection{PRS für ganzzahlige Rechnung
\label{resultante:subsection:ganzzahlig}}
Angenommen, die Terme eines Polynoms \(P\) sind in absteigender Potenz
sortiert, dann ist der Koeffizient des ersten Terms der Leitkoeffizient
\(\lc(P)\).
\index{Leitkoeffizient}%
\index{lc(P)@$\operatorname{lc}(P)$}%
Daher kann man die normale Division und deren Rest durch eine neue
Pseudo-Division mit Pseudo-Rest ersetzen, in der sich die Resultate
\index{Pseudo-Division}%
\index{Pseudo-Rest}%
der normalen Divison und der normale Rest nur um bestimmte Faktoren
unterscheiden.
Die Pseudo-Division hat rein ganzzahlige Pseudo-Quotienten
\(Q=\pquo(f,g)\) mit ganzzahligem Pseudo-Rest \(R=\prem(f,g)\).
Anders als in der gewöhlichen Polynom-Division mit Rest entstehen dabei
keine Brüche, da mit einer Potenz des Leitkoeffizienten des Divisors
multipliziert wird, sodass
\begin{equation}
    b_i^{\delta_i+1}P_i-Q_iP_{i+1}=R_i \quad \text{und} \quad \deg(R)<\deg(P_{i+1})\quad \text{für} \quad i=1,\dots,k-1 \, \text{ist},
    \label{resultante:equation:algpdivi}
\end{equation}
wobei \(b_i\) der Leitkoeffizient von \(P_{i+1}\) (\(\lc(P_{i+1})\)),
\(\delta_i=\deg(P_i)-\deg(P_{i+1})\) und \(P_{i+2}=R_i\) ist.
Bei Anwendung derselben Vorgehensweise auf \(P_2\) und \(P_3\), und
danach auf \(P_3\) und \(P_4\) usw., entsteht eine Folge
\begin{equation}
    P_1,P_2,P_3,\dots,P_k,P_{k+1}=0.
    \label{resultante:equation:prs}
\end{equation}
Eine solche Folge heisst Polynom-Restfolge (PRS).
\index{Polynom-Restfolge}%
\index{polynomial remainder sequence}%
\index{PRS}%

\begin{beispiel}
    Aus den Polynomen \(a(X)=4X^4-7X^3+22X^2-24X+5\) und
    \(b(X)=8X^3+8X^2-16X\) soll die PRS gebildet werden.
    Begonnen wird mit \(P_1=a(X)\) und \(P_2=b(X)\), dabei ist
    \begin{equation*}
        b_1 = 8, \qquad \delta_1=4-3=1.
    \end{equation*}
    Aus (\ref{resultante:equation:algpdivi}) folgt im ersten Schritt
    \begin{equation}
        8^{1+1}P_1-Q_1P_2=R_1,
    \end{equation}
    nach der Pseudo-Division erhält man \(Q_1=32X-88\) und
    \(R_1=P_3=2624X^2-2944X+320\).
    Der nächste Schritt hat
    \begin{equation*}
        b_2 = 2624, \qquad \delta_2=2-1=1,
    \end{equation*}
    daraus folgt die neue Division
    \begin{equation}
        2624^{1+1}P_2-Q_2P_3=R_2,
    \end{equation}
    was \(Q_2=20992X+44544\) und \(R_2=P_4=14254080X-14254080\) ergibt.
    Da \(R_2\neq0\) ist, geht der Algorithmus weiter mit
    \begin{equation*}
        b_3 = 14254080, \qquad \delta_3=1-1=0,
    \end{equation*}
    und
    \begin{equation}
        14254080^{1}P_3-Q_3P_4=R_3,
    \end{equation}
    daraus resultiert \(Q_3=37402705920X-4561305600\) und \(R_3=P_5=0\).
    Die aus den Rest-Polynomen entstehende Folge ist
    \begin{align*}
            P_1&=4X^4-7X^3+22X^2-24X+5, \\
            P_2&=8X^3+8X^2-16X, \\
            P_3&=2624X^2-2944X+320, \\
            P_4&=14254080X-14254080, \\
            P_5&=0.
	\qedhere
    \end{align*}
    \label{resultante:beispiel:prs}
\end{beispiel}
Wie man in Beispiel \ref{resultante:beispiel:prs} beobachten kann,
wird Ziel a) erreicht, aber b) völlig verfehlt.
Die Koeffizienten der PRS explodieren.
Ist das letzte Polynom der Folge \(P_{k+1}=0\), dann ist \(P_k\)
der unnormierte \(\ggt\), in Beispiel \ref{resultante:beispiel:prs}
\(P_4=14254080(X-1)\).


\subsection{Subresultanten-PRS
\label{resultante:subsection:subresprs}}
Um das Koeffizienten-Wachstum in den Griff zu bekommen, wird die Folge
mit \(\beta_i\) normiert.
Dadurch wird
\begin{equation}
    \beta_i P_i=\prem(P_{i-2},P_{i-1})\quad \text{für} \quad i=3,\dots,k+1.
    \label{resultante:equation:prsnorm}
\end{equation}
Nun muss nur noch ein geeigneter Faktor \(\beta_i\) gefunden werden,
dieser kann im Prinzip beliebig gewählt werden.
Je nach Wahl des \(\beta_i\) ist das Koeffizienten-Wachstum besser oder
schlechter kontrollierbar, oder es kann sichergestellt werden, dass die
Koeffizienten immer in \(\mathbb{Z}\) bleiben.

Der künstliche Faktor, der in der Pseudo-Division vorkommt, ist bereits
bekannt und kommt als Skalierungsfaktor \(\beta_i\) infrage.
Es ist nicht ganz einfach, einen solchen \(\beta_i\)-Faktor zu finden,
der das Ziel a) immer noch erreicht.
Dieser \(\beta_i\)-Faktor kann in der PRS wie folgt festgelegt werden.
Um eine PRS zu formulieren, sei \(p_i=\lc(P_i)\) und man wählt
\begin{equation}
    \begin{aligned}
        \beta_3&=(-1)^{\delta_1+1}\\
        \beta_i&=(-1)^{\delta_{i-2}+1}p_{i-2}h^{\delta_{i-2}}_{i-2}, \quad i=4,\dots,k+1,
    \end{aligned}
    \label{resultante:equation:beta}
\end{equation}
wobei die Hilfsfaktoren \(h_i\)
\begin{equation}
    \begin{aligned}
        h_2&=p_2^{\delta_1}\\
        h_i&=p_i^{\delta_{i-1}}h^{1-\delta_{i-1}}_{i-1}, \quad i=3,\dots,k
    \end{aligned}
    \label{resultante:equation:h}
\end{equation}
sind.
Das ist die Brown-Normierung.
\index{Brown-Normierung}%
W.~S.~Brown hat in \cite{resultant:PRSalg} die Subresultanten-PRS
\index{Brown, W.~S.}%
weiterentwickelt, die George E.~Collins bereits in
\index{Collins, George E.}%
\cite{resultant:Collins} beschreibt.

Daraus folgt die Subresultanten-PRS.
Die daraus konstruierte Subresultanten-PRS mit \(G_1=P_1,\,G_2=P_2,\,G_3,\dots,\,G_k\), wird mit
\begin{equation}
    \begin{aligned}
        G_3&=(-1)^{\delta_1+1}\prem(G_1,G_2),\\
        G_i&=\frac{(-1)^{\delta_{i-2}+1}\prem(G_{i-2},G_{i-1})}{g_{i-2}h^{\delta_{i-2}}_{i-2}},\quad i=4,\dots,k+1
    \end{aligned}
    \label{resultante:equation:algsubprs}
\end{equation}
berechnet.
Die Iteration stoppt, sobald \(\prem(G_{k-1},G_k)=0\) ist.
Dann ist
\begin{equation}
    \ggt(P_1,P_2)=\pp(G_k), \qquad \Res(P_1,P_2)=0,
    \label{resultante:equation:ggtres}
\end{equation}
wobei \(\pp\) das primitive Polynom ist, im Beispiel
\ref{resultante:beispiel:prs} ist \(\pp(P_4)=X-1\).
Dabei haben die Koeffizienten den \(\ggt=1\).
Der Algorithmus soll mit folgendem Beispiel verdeutlicht werden.

\begin{beispiel}
    Zunächst wird mit \(G_1=P_1=a(X)=4X^4-7X^3+22X^2-24X+5\) und \(G_2=P_2=b(X)=8X^3+8X^2-16X\) begonnen.
    Für die Bestimmung von \(G_3\) wird nach (\ref{resultante:equation:algsubprs}) mit
    \begin{equation*}
        G_3=(-1)^{\delta_1+1}\prem(G_1,G_2)
    \end{equation*}
    begonnen, dabei ist
    \begin{equation*}
        \delta_1=\deg(G_1)-\deg(G_2)=1.
    \end{equation*}
    Daraus folgt
    \begin{equation*}
        G_3=\prem(G_1,G_2)=2624X^2-2944X+320.
    \end{equation*}
    Als nächstes wird \(G_4\) mit
    \begin{equation*}
        G_4=\frac{(-1)^{\delta_2+1}\prem(G_2,G_3)}{g_2h^{\delta_2}_2}
    \end{equation*}
    bestimmt, dafür braucht es
    \begin{equation*}
        \delta_2=\deg(G_2)-\deg(G_3)=1, \quad g_2=8, \quad h^{\delta_2}_2=g^{\delta_2}_2=8.
    \end{equation*}
    Das in die Gleichung für \(G_4\) eingesetzt ergibt
    \begin{equation*}
        G_4=\frac{\prem(G_2,G_3)}{8\cdot8}=222720X-222720.
    \end{equation*}
    Da \(G_4\neq0\) ist, geht der Algorithmus weiter mit
    \begin{equation*}
        G_5=\frac{(-1)^{\delta_3+1}\prem(G_3,G_4)}{g_3h^{\delta_3}_3}.
    \end{equation*}
    Hier müssen \(\delta_3, g_3, h_3\) nicht mehr bestimmt werden, da bereits \(\prem(G_3,G_4)=0\) ist und somit \(G_5=0\).
    Daraus entsteht die Subresultanten-PRS
    \begin{equation*}
        \begin{aligned}
            G_1&=4X^4-7X^3+22X^2-24X+5, \\
            G_2&=8X^3+8X^2-16X, \\
            G_3&=2624X^2-2944X+320, \\
            G_4&=222720X-222720, \\
            G_5&=0.
        \end{aligned}
    \end{equation*}
    Der \(\ggt\) ist nach (\ref{resultante:equation:ggtres}) links \(\ggt(P_1,P_2)=\pp(G_4)=X-1\).
    \label{resultante:beispiel:subprs}
\end{beispiel}
Wenn man \(P_4\) aus der gewöhlichen PRS mit \(G_4\) aus der
Subresultanten-PRS vergleicht, ist ganz deutlich zu sehen, welchen
Einfluss die Brown-Normierung auf die Koeffizienten hat.
Die Koeffizienten in \(G_4\) sind deutlich kleiner als die von \(P_4\).
Ziel a) wird immer noch erreicht und man ist dem Ziel b) deutlich näher
gekommen.

\subsection{Anwendung der Subresultanten-PRS
\label{resultante:subsection:subprs}}
\begin{beispiel}
    Wie in Beispiel \ref{resultante:beispiel:gutesIntegral} soll die
    Stammfunktion von
    \begin{equation*}
        \int\frac{1}{z^3+z}\,dz
    \end{equation*}
    bestimmt werden.
    Dafür muss die Resultante (\ref{resultante:equation:resZgut})
    \begin{equation}
        R(c)=\Res_z\left(-3cz^2+1-c,z^3+z\right)
        \label{resultante:equation:resprs}
    \end{equation}
    berechnet werden. 
    Dies soll mit der Subresultanten-PRS durchgeführt werden.
    Damit man den Überblick über die einzelnen Faktoren der Berechnung
    behält, werden sie in einer Tabelle notiert.
    Die \(G_i\) werden nach (\ref{resultante:equation:algsubprs}) in
    \begin{equation*}
        \begin{NiceArray}{|c|l|c|c|c|c|}[cell-space-limits=3pt]
            \hline
            i & G_i & \delta_i & n_i & g_i & h_i \\
            \hline
            1 & z^3+z & 1 & 3 & 1 & - \\
            2 & -3cz^2+1-c & 1 & 2 & -3c & -3c \\
            3 & 3c(2c+1)z & 1 & 1
            & 9c^2(2c+1)^2
            & -27c^3(2c+1)^2 \\
            4 & (1-c)(2c+1)^2 & - & 0 & - & - \\
            \hline
        \end{NiceArray}
        \label{resultante:equation:tabsubprs}
    \end{equation*}
    berechnet.
    Dabei ist \(n_i\) der Grad von \(G_i\), \(g_i\) der Leitkoeffizient
    \(\lc(G_i)\) und \(h_i\) der Hilfsfaktor.
    Da \(G_4\) konstant bezüglich \(z\) ist, ist \(G_4\) die Resultante
    \(\Res_z\left(-3cz^2+1-c,z^3+z\right)\).
    Die Resultante soll \(\Res_z=0\) sein, das wird erreicht, wenn
    \(c_1=1\) oder \(c_2=-\frac{1}{2}\) ist.
    Nun kann der \(\ggt\) mit \(G_3\) der darüber liegenden Zeile
    bestimmt werden, für \(c_1=1\):
    \begin{equation}
        \begin{aligned}
            G_3(c_1)&=3c_1(2c_1+1)z=3\cdot1(2\cdot1+1)z=9z\\
            \Rightarrow\quad \ggt&=\pp\left(G_3(c_1)=9z\right)=z.
        \end{aligned}
    \end{equation}
    Analog gilt für \(c_2=-\frac{1}{2}\):
    \begin{equation}
        G_3(c_2)
        =
        3c_2(2c_2+1)z
        =
        3\cdot-\frac{1}{2}\biggl(2\cdot-\frac{1}{2}+1\biggl)z=0,
    \end{equation}
    Hier geschieht Besonderes: \(G_3\) wird ebenfalls Null, wenn \(c_2\)
    eingesetzt wird.
    Um den \(\ggt\) zu bestimmen, muss daher \(G_2\) in der
    darüberliegenden Zeile verwendet werden:
    \begin{equation}
        \begin{aligned}
            G_2(c_2)&=-3c_2z^2+1-c_2=-3\cdot-\frac{1}{2}z^2+1-\biggl(-\frac{1}{2}\biggr)=1.5z^2+1.5\\
            \Rightarrow\quad \ggt&=\pp\left(G_2(c_2)=1.5z^2+1.5\right)=z^2+1.
        \end{aligned}
    \end{equation}
    Setzt man wieder die gefundenen Nullstellen mit den dazugehörigen
    \(\ggt\) zusammen, erhält man die Stammfunktion:
    \begin{equation*}
        \int\frac{1}{z^3+z}\,dz=c_1\log v_1(z)+c_2\log v_2(z)=1\log z-\frac{1}{2}\log(z^2+1). \qedhere
    \end{equation*}
    \label{resultante:beispiel:subprsRT}
\end{beispiel}
Das Beispiel \ref{resultante:beispiel:subprsRT} zeigt einen
Algorithmus, mit dem dasselbe erreicht wird, wie in Beispiel
\ref{resultante:beispiel:gutesIntegral}: einen Weg, die Resultante zu
berechnen und die Nullstelle zweier Polynome zu finden.
